\chapter{Граница происхождения общего дираковского подъёма}
\label{ch:v8-full-field-a4-dirac-lift-origin-gate}

Пусть конечный оператор и полный внутренний репер фиксированы. Для любого
$r>0$ допустим четырёхмерный оператор произведения
\begin{equation}
 \mathcal D_r=rD_M\otimes I_{21}+\gamma^5\otimes D_F.
 \label{eq:v8-product-dirac-scale-family}
\end{equation}
Для каждого внутреннего наблюдаемого $a$ выполнено
\begin{equation}
 [\mathcal D_r,I_4\otimes a]=\gamma^5\otimes[D_F,a],
 \label{eq:v8-internal-calculus-scale-invariance}
\end{equation}
то есть все сорок два направления конечного дифференциального исчисления не
видят множитель $r$. В то же время главный внешний символ меняется:
\begin{equation}
 \sigma_1(D_M)^2=I_4,
 \qquad \sigma_1(2D_M)^2=4I_4.
 \label{eq:v8-external-symbol-scale-inequivalence}
\end{equation}

\begin{theorem}[Устойчивость относительного веса]
Клиффордово отношение $3:1$, найденное в общем $a_4$-блоке, одинаково для
всего семейства $\mathcal D_r$: внешний масштаб умножает общий локальный
блок, но не меняет относительный спинорный след.
\end{theorem}

\begin{nogocriterion}[Конечный родитель не выбирает геометрию базы]
Совпадение внутренней алгебры, конечного оператора, суперсвязности и всех
сорока двух внутренних коммутаторов не фиксирует $r$, метрику базы или
спектральный функционал. Поэтому общий $a_4$-подъём является совместимой
архитектурой, но не безусловным следствием конечного родителя.
\end{nogocriterion}

\begin{protocol}
\begin{enumerate}
 \item размерность спинорного репера: \textbf{$4$};
 \item внутренних направлений: \textbf{$42$};
 \item дефект внутренних коммутаторов при $r:1\mapsto2$: \textbf{$0$};
 \item квадрат внешнего символа: \textbf{$1\mapsto4$};
 \item отношение $a_4$-весов: \textbf{$3:1$, сохраняется};
 \item внешний геометрический масштаб: \textbf{не выбран};
 \item следующий гейт: \textbf{селектор геометрии пространственно-временной базы}.
\end{enumerate}
\end{protocol}

Машинный сертификат сохранён в
\path{s2t_v8_full_field_a4_dirac_lift_origin_gate_results.json}.